\section{The actual theta kernel and the selected-prime column estimate}
\label{tvc-section}

The moment and incidence hypotheses in Section~\ref{cv-section}
have now been proved for the normalized theta functions of the
construction. We make the remaining identification with the
polynomial kernel explicit. In particular, an arbitrary projection
onto the target representation has no effect, and ordinary mixed
derivatives introduce no factorial denominator. We then state exactly
which column estimate follows. This concerns the selected primes;
the estimates at the auxiliary level primes are separate.
\cite{alpo-uniform-source}

\subsection{From mixed derivatives to the polynomial map}

Retain the ranks $a<b$, embedding set $I$, dominant weights
$\lambda,\mu$, and nonnegative determinant exponents $n_\tau$
of \eqref{cv-weight-compatibility}. For the source variables
$X_{\tau,\alpha s}$, target variables $Y_{\tau,\beta j}$,
and new variables $Z_{\tau,\alpha\beta}$, put
\[
 \begin{split}
 F^\sharp(X)&=F(X)\prod_{\tau\in I}\det(X_\tau)^{n_\tau},\\
 \mathcal Q_F(Y,Z)
   &=F^\sharp\bigl((Z_\tau Y_{\tau,[1,a]})_{\tau\in I}\bigr)
     =\sum_m Z^m Q_m(F)(Y).
 \end{split}
                                                               \tag{TC1}
\]
Here $Y_{\tau,[1,a]}$ denotes the first $a$ columns of $Y_\tau$.
This is the universal polynomial of
\eqref{cv-universal-expansion}. The order of its two new matrix
indices is significant: $Z$ is an $a$ by $b$ matrix, whereas
the analytic theta variable is a $b$ by $a$ matrix. Thus
$Z_{\tau,\alpha\beta}$ corresponds to differentiation in
the analytic coordinate $(\tau,\beta,\alpha)$.

Let $G$ be holomorphic near zero in those analytic coordinates,
and let $c_\tau\ne0$ be the fixed direction constants. Define
\[
 J_m(G)=
  \left(\prod_{\tau,\alpha,\beta}
                  c_\tau^{m_{\tau,\alpha\beta}}\right)
  \left(\prod_{\tau,\alpha,\beta}
   \partial_{z_{\tau,\beta,\alpha}}^{m_{\tau,\alpha\beta}}
          G\right)(0).
                                                               \tag{TC2}
\]
Holomorphic mixed derivatives commute. Their multilinear definition
also shows that the result is independent of the order in which the
repeated directions are listed. In particular $J_m(G)$ is the
derivative, and is $m!$ times the corresponding Taylor coefficient
when all $c_\tau=1$.

\begin{lemma}\label{tvc-jet-map}
For any scalar family $(J_m)$, the operation on polynomials that
first multiplies by $\prod_\tau\det(X_\tau)^{n_\tau}$,
then substitutes
\[
              X_{\tau,\alpha s}
                   \longmapsto
                  \sum_\beta Y_{\tau,\beta s}Z_{\tau,\alpha\beta},
\]
and finally sends each monomial $Y^vZ^m$ to $J_mY^v$, is
precisely
\[
                       F\longmapsto\sum_m J_mQ_m(F).
                                                               \tag{TC3}
\]
Its image lies in $V_\mu$ for $F\in V_\lambda$. Consequently
composing it with any linear projection of the full target polynomial
ring onto $V_\mu$ which is the identity on $V_\mu$ leaves it unchanged.
\end{lemma}
\begin{proof}
After the first two operations the polynomial is exactly (TC1).
Applying the stated monomial rule gives (TC3) by linearity.
Lemma~\ref{cv-substitution} proves that $\mathcal Q_F(Y,Z)$,
as a polynomial in $Y$, satisfies the defining column degrees and
column-addition invariances of $V_\mu$. Comparing coefficients of
the independent variables $Z$ shows that each $Q_m(F)$ satisfies
these same identities. Hence every finite linear combination in
(TC3) belongs to $V_\mu$. A projection equal to the identity there
therefore fixes it. This argument makes no choice of a complement
and uses no integrality of a projection.
\end{proof}

For the actual lattices $N,N'$ and finite selected set $S$, put
\[
 H=\{x\in\operatorname{Hom}_E(E^a,E^b):xN_P\subseteq N'_P
                                                \ (P\in S)\}.
\]
For positive integers $c_A,c_B$, choose representatives $x$ of
$c_A^{-1}H/c_BH$. This quotient is finite by
Lemma~\ref{jcf-finite-indicators}. Let $G_x$ be its normalized
theta indicator, including division by the determinant-character
factor as in Section~\ref{jcf-section}. Let $\Xi(x)$ be the actual
finite test function, and let $\eta(N,N')$ be the ratio of the
remaining local frame and determinant characters. Set
\[
                 J_m=\sum_{x\in c_A^{-1}H/c_BH}
                                  \Xi(x)J_m(G_x).
                                                               \tag{TC4}
\]
All the quantities in this formula belong to the one field $K$
of Proposition~\ref{jcf-common-field}, after the fixed finite
choices have been made. Only the finitely many $m$ of
\eqref{cv-finite-degrees} are used.

\begin{proposition}\label{tvc-actual-kernel}
The actual polynomial theta kernel for this pair of decorated
lattices is
\[
                  K(N,N')F=\eta(N,N')\sum_m J_mQ_m(F).
                                                               \tag{TC5}
\]
Here $J_m$ is exactly (TC4), with the derivatives (TC2).
There is no additional factorial, lattice index, or projection
factor in (TC5).
\end{proposition}
\begin{proof}
For a fixed representative $x$, the polynomial theta operation is
the determinant multiplication and substitution of Lemma~\ref{tvc-jet-map},
followed by the derivative functional $Z^m\mapsto J_m(G_x)$.
That lemma identifies it with $\sum_mJ_m(G_x)Q_m$ and removes
any projection onto the target representation. Multiply by
$\Xi(x)$ and sum. Both sums are finite, so interchanging them
gives (TC4)--(TC5). The factor $\eta(N,N')$ is independent of $x$
and can be taken outside the sum.

The finite indicator $G_x$ already contains the Fourier average
and the positive square-root normalization proved in
Section~\ref{csp-section}. The present operation neither inserts
nor removes them. Likewise (TC2) is the ordinary repeated
derivative used there, not a divided Taylor coefficient.
This proves the normalization assertion. The admissibility of
the representatives is also exact: $c_Ax\in H$ implies
$c_Ac_Bx\in H$, so each term lies in the permitted level
$c_Ac_B$ and the indicator is evaluated at level $c_B$.
All equalities descend from the complex embedding to $K$ by
its injectivity.
\end{proof}

\subsection{The resulting local divisibility}

We retain all the geometric and coefficient data of
Section~\ref{cv-section}. Thus $[E:\mathbb Q]=2f$,
$\pi\bar\pi=p$, the $f$ primes dividing $\pi$ have norm $p$,
and the ranks are positive even integers
\[
                 N_{i+1}-N_i=t+1=2s_0\quad(0\le i<m).
\]
The CM elliptic ideal and the inverse half different are units at
$p$, and $p\ge5$. The source and target weights satisfy
\eqref{cv-weight-compatibility} at each transition. Choose
nonnegative integers $W_\tau$ with \eqref{cv-weight-bounds}, and suppose the
target hook bound \textup{(CV11)} holds at every transition and
\[
                              p>2fN_m+1.                       \tag{TC6}
\]
All other local levels in the test function are prime to $p$,
with integral values at the coefficient place in question, as in
the construction of Section~\ref{csp-section}. These conditions
are retained because they are used in the moment and character
unit calculations.

Let $\mathscr S_i$ be the local coefficient lattice of equivariant
class sections at rank $N_i$. Let $\mathcal T$ be the distinguished
set of embeddings, $r=|\mathcal T|$, and let $U_i$ be the product
of the upper-neighbor sums at its corresponding selected primes.
Let $K_i$ be the actual class-sum theta transition and put
$C_m=K_{m-1}\cdots K_0$.

\begin{theorem}\label{tvc-column}
At a coefficient place $w\mid p$ of selected type, put
\[
 e=2\sum_\tau W_\tau+fN_m,\qquad
 b_m=me+s_0r\frac{m(m+1)}2.
                                                               \tag{TC7}
\]
Then
\[
                  p^{b_m}C_mU_0^m\mathscr S_0
                                          \subseteq\mathscr S_m.
                                                               \tag{TC8}
\]
If the initial lattice is $\mathscr N=p^AU_0^m\mathscr S_0$
with $A\ge0$, and $v_w(c)\ge v_w(p)b_m$, then
\[
                              cC_m\mathscr N\subseteq\mathscr S_m.
                                                               \tag{TC9}
\]
At a coefficient place of conjugate type the column preserves
the original integral class lattice without a power of $p$.
Applying any of the integral lowering maps of
Section~\ref{atl-section}, and then pairing with an integral row,
introduces no further denominator at $w$.
\end{theorem}
\begin{proof}
Proposition~\ref{tvc-actual-kernel} identifies the actual operator
with the operator in \textup{(CV2)}. The moment laws of
Section~\ref{gm-section} and the two support calculations of
Section~\ref{csp-section} supply exactly \textup{(CV9)--(CV10)}:
the selected support has scalar $p^{-fN_i}$ and the conjugate
support has scalar one. Lemma~\ref{csp-coordinate-units} gives
the unit direction factors. The approximations use one integral
finite sum for all the required moments. Thus they are the
simultaneous approximations in
\eqref{cv-moment-approximation}, not merely separate bounds on
individual derivatives.

The character ratio is a unit by
Lemma~\ref{cv-character-unit}, and the hook bound identifies the
target Weyl and dual-Weyl lattices. Condition (TC6) makes every
class stabilizer order a unit by
Lemma~\ref{cv-stabilizer-unit}. Theorem~\ref{cv-two-kernels}
and the residue-removal argument therefore give
\[
                      p^eU_{i+1}K_i\mathscr S_i
                                                \subseteq\mathscr S_{i+1}.
\]
Theorem~\ref{hec-global-identity}, at each of the $r$ distinct
primes, gives
\[
                            U_{i+1}K_i=p^{s_0r}K_iU_i.
\]
The different local neighbor operators commute since they change
different local components. All these identities hold over the
same field $K$ by Proposition~\ref{jcf-common-field}.
Proposition~\ref{cv-column} now proves (TC8)--(TC9), with the
integer exponent (TC7). At conjugate-type places the second
bound of Theorem~\ref{cv-two-kernels} already preserves each
original integral lattice, so their composition does too.

The final lowering maps preserve the local lattices by
Section~\ref{atl-section}. The row pairing of
Section~\ref{cp-section} has only the stated class stabilizer
denominators, which are units by (TC6). This proves the last
assertion.
\end{proof}

The exponent (TC7) is the exponent $mE_p$ of
\eqref{cv-final-exponent}. The selected-prime part of the column
valuation argument is therefore complete under the displayed
construction hypotheses. This does not prove that
$U_0^m\mathscr S_0$ or its theta image is nonzero, does not supply
the estimates at the two auxiliary level primes, and does not
replace the archimedean norm and period comparison needed for
the final arithmetic inequality. Those are separate mathematical
assertions.


% Scalar rigidity for the actual Hermitian flags.
% Labels in this fragment use the prefix frg-.
